\documentclass[a4paper,11pt]{article}
\usepackage[T2A]{fontenc}
\usepackage[utf8]{inputenc}
\usepackage[russian]{babel}
\usepackage{amsmath,amssymb,mathtools}
\usepackage{geometry}
\geometry{left=23mm,right=23mm,top=21mm,bottom=20mm}
\usepackage{xcolor}
\definecolor{accent}{HTML}{087F80}
\definecolor{darkaccent}{HTML}{07555A}
\definecolor{textgray}{HTML}{4F5559}
\usepackage{tikz}
\usetikzlibrary{arrows.meta,positioning,shapes.geometric,calc}
\usepackage{booktabs,tabularx,array}
\usepackage{enumitem}
\usepackage{fancyhdr}
\usepackage{microtype}
\usepackage{hyperref}
\hypersetup{colorlinks=true,linkcolor=darkaccent,urlcolor=darkaccent}
\usepackage{setspace}
\setstretch{1.08}
\setlength{\parindent}{0pt}
\setlength{\parskip}{3pt}
\AtBeginDocument{\setlength{\abovedisplayskip}{7pt}\setlength{\belowdisplayskip}{7pt}\setlength{\abovedisplayshortskip}{4pt}\setlength{\belowdisplayshortskip}{4pt}}
\newcommand{\secname}[1]{\vspace{5pt}{\large\color{darkaccent}#1}\par\vspace{2pt}\noindent{\color{accent}\rule{\linewidth}{.5pt}}\vspace{3pt}}
\newcommand{\ex}[1]{\par\smallskip{\color{darkaccent}\textbf{Пример.}} #1\par}
\newcommand{\lead}[1]{{\color{darkaccent}\textbf{#1}}}
\pagestyle{fancy}
\fancyhf{}
\fancyhead[L]{\footnotesize\color{textgray}Тригонометрические преобразования}
\fancyhead[R]{\footnotesize\color{textgray}Марк Гукин}
\renewcommand{\headrulewidth}{.35pt}
\renewcommand{\headrule}{\hbox to\headwidth{\color{accent}\leaders\hrule height \headrulewidth\hfill}}
\fancyfoot[C]{\small\thepage}
\setlength{\headheight}{23pt}
\begin{document}
\thispagestyle{empty}
\vspace*{27mm}
\begin{center}
{\large\color{darkaccent} МАТЕМАТИКА}\par
\vspace{20mm}
{\fontsize{24}{30}\selectfont Чек-лист}\par
\vspace{12mm}
{\fontsize{19}{24}\selectfont\color{accent} Тригонометрические преобразования\\[3mm]
Связанные углы и двойной угол}\par
\vspace{16mm}
{\large Подготовка к ЕГЭ по профильной математике}\par
\vspace{36mm}
{\large Марк Гукин}\par
\vspace{12mm}
{\large 1 ноября 2026 г.}\par
\vfill
{\small Лёвин Артём Александрович\\[2mm] эксклюзивно для Марка Гукина}\par
\vspace{17mm}
\begin{tikzpicture}
\draw[accent,line width=.85pt] (-4,0) .. controls (-2.55,.72) and (-1.15,-.55) .. (0,0)
 .. controls (1.1,.58) and (2.35,-.7) .. (4,0);
\fill[accent] (0,0) circle (1.6pt);
\end{tikzpicture}
\end{center}
\clearpage

\secname{1. Главная идея: привести углы и функции к одному виду}

В вычислительных тригонометрических задачах ищите связь между углами. Особое внимание уделяйте случаям, когда углы \lead{в сумме дают $90^\circ$}, \lead{отличаются в два раза} или отличаются на целое число периодов.

\begin{center}
\begin{tikzpicture}[>=Latex, every node/.style={font=\small}]
\node (a) {Сравнить углы};
\node[right=12mm of a] (b) {Выбрать формулу};
\node[right=12mm of b] (c) {Получить одинаковые углы и функции};
\draw[->,accent] (a)--(b);
\draw[->,accent] (b)--(c);
\end{tikzpicture}
\end{center}

\secname{2. Дополнительные углы и основное тождество}

Если $\alpha+\beta=90^\circ$, то
\[
\cos(90^\circ-\alpha)=\sin\alpha,\qquad
\sin(90^\circ-\alpha)=\cos\alpha.
\]
Эти формулы следуют из формул синуса и косинуса разности с учётом $\sin90^\circ=1$ и $\cos90^\circ=0$.

\ex{На занятии вычислялось выражение
\[
\frac{51\cos4^\circ}{\sin86^\circ}+8
=\frac{51\sin86^\circ}{\sin86^\circ}+8=\boxed{59}.
\]
Здесь $4^\circ+86^\circ=90^\circ$.}

Для острых дополнительных углов:
\[
\tg(90^\circ-\alpha)=\frac{\cos\alpha}{\sin\alpha}=\ctg\alpha.
\]
\ex{\[
46\tg7^\circ\tg83^\circ
=46\frac{\cos83^\circ}{\sin83^\circ}
\frac{\sin83^\circ}{\cos83^\circ}=\boxed{46}.
\]}
\lead{Ограничение.} $\tg x=\dfrac{\sin x}{\cos x}$ определён лишь при $\cos x\ne0$. В частности, $\tg90^\circ$ не существует. Подставлять его в формулу тангенса разности нельзя.

\secname{3. Основное тригонометрическое тождество}
\[
\sin^2\alpha+\cos^2\alpha=1.
\]
\ex{Если углы дополняют друг друга до $90^\circ$, можно получить квадрат синуса и квадрат косинуса одного угла:
\[
\frac{19}{\cos^237^\circ+1+\cos^253^\circ}
=\frac{19}{\sin^253^\circ+1+\cos^253^\circ}
=\boxed{\frac{19}{2}=9{,}5}.
\]}

\secname{4. Формулы двойного угла: используем в обе стороны}
\begin{align*}
\sin2\alpha&=2\sin\alpha\cos\alpha,\\
\cos2\alpha&=\cos^2\alpha-\sin^2\alpha,\\
            &=2\cos^2\alpha-1,\\
            &=1-2\sin^2\alpha.
\end{align*}
Важно уметь \lead{раскрывать} формулы и \lead{сворачивать} выражения обратно, когда нужная комбинация уже дана.

\ex{Одновременно замечаем удвоенный угол и дополнительные углы:
\begin{align*}
\frac{5\sin98^\circ}{\sin49^\circ\sin41^\circ}
&=\frac{5\cdot 2\sin49^\circ\cos49^\circ}
{\sin49^\circ\sin41^\circ}\\
&=\frac{10\cos49^\circ}{\sin41^\circ}
=\frac{10\sin41^\circ}{\sin41^\circ}=\boxed{10}.
\end{align*}}

\secname{5. Противоположные знаки: вынести минус}
\[
\sin^2\alpha-\cos^2\alpha
=-\bigl(\cos^2\alpha-\sin^2\alpha\bigr)
=-\cos2\alpha.
\]
\ex{\[
\frac{24(\sin^217^\circ-\cos^217^\circ)}{\cos34^\circ}
=\frac{-24(\cos^217^\circ-\sin^217^\circ)}{\cos34^\circ}
=\frac{-24\cos34^\circ}{\cos34^\circ}
=\boxed{-24}.
\]}
\lead{Проверка.} Сокращать разрешено только \emph{одинаковые ненулевые множители}. Если скобки отличаются знаками, сначала вынесите минус. Для примера $\cos34^\circ\neq0$.

\secname{6. Общий множитель и свёртывание двойного угла}

Если в выражении есть общий множитель, его удобно вынести за скобки. Затем проверьте, можно ли свернуть оставшуюся разность или сумму квадратов.

\ex{\begin{align*}
\sqrt{50}\cos^2\frac{9\pi}{8}-\sqrt{50}\sin^2\frac{9\pi}{8}
&=\sqrt{50}\left(\cos^2\frac{9\pi}{8}-\sin^2\frac{9\pi}{8}\right)\\
&=\sqrt{50}\cos\frac{9\pi}{4}
=\sqrt{50}\cos\left(2\pi+\frac\pi4\right)\\
&=\sqrt{50}\frac{\sqrt2}{2}=\boxed{5}.
\end{align*}}

Для второй формы удобно выделить общий множитель так, чтобы в скобках получилось $2\cos^2\alpha-1$.
\ex{\begin{align*}
4\sqrt2\cos^2\frac{15\pi}{8}-2\sqrt2
&=2\sqrt2\left(2\cos^2\frac{15\pi}{8}-1\right)\\
&=2\sqrt2\cos\frac{15\pi}{4}
=2\sqrt2\cos\left(4\pi-\frac\pi4\right)=\boxed{2}.
\end{align*}}

\secname{7. Большие и отрицательные углы}

Для целого $k$ верны периодичность и чётность/нечётность:
\begin{align*}
\sin(\alpha+2\pi k)&=\sin\alpha,& \cos(\alpha+2\pi k)&=\cos\alpha,\\
\sin(-\alpha)&=-\sin\alpha,& \cos(-\alpha)&=\cos\alpha,\\
\tg(-\alpha)&=-\tg\alpha,& \ctg(-\alpha)&=-\ctg\alpha.
\end{align*}
В градусах период синуса и косинуса равен $360^\circ$.

\ex{\[
\sin\left(-\frac{27\pi}{4}\right)
=-\sin\left(6\pi+\frac{3\pi}{4}\right)
=-\sin\frac{3\pi}{4}=\boxed{-\frac{\sqrt2}{2}}.
\]}

\lead{Градусы и радианы:} $\alpha_{\text{рад}}=\alpha_{\text{град}}\cdot\dfrac{\pi}{180}$, поэтому
\[
135^\circ=\frac{135\pi}{180}=\frac{3\pi}{4}=\pi-\frac\pi4.
\]
\ex{\[
-18\sqrt2\sin(-135^\circ)
=18\sqrt2\sin135^\circ
=18\sqrt2\cdot\frac{\sqrt2}{2}=\boxed{18}.
\]}

\secname{8. Знаки в формулах суммы и разности}
\[
\begin{aligned}
\cos(90^\circ+\alpha)&=-\sin\alpha,\\
\cos(90^\circ-\alpha)&=\sin\alpha,\\
\sin(360^\circ-\alpha)&=-\sin\alpha.
\end{aligned}
\]
\lead{Ловушка со знаком:} $\cos105^\circ=-\sin15^\circ$. Отсюда
\[
-8\cos15^\circ\cos105^\circ
=8\cos15^\circ\sin15^\circ
=4\sin30^\circ=\boxed{2}.
\]

\secname{9. Итоговый пример с несколькими преобразованиями}
\begin{align*}
\frac{6\sin95^\circ\cos95^\circ}{\sin170^\circ}
&=\frac{3\sin190^\circ}{\sin170^\circ}
=\frac{3\sin(360^\circ-170^\circ)}{\sin170^\circ}\\
&=\frac{-3\sin170^\circ}{\sin170^\circ}=\boxed{-3}.
\end{align*}
\lead{Сигналы:} в числителе есть произведение синуса и косинуса одного угла; после формулы двойного угла появляются углы, дополняющие друг друга до $360^\circ$.

\clearpage
\thispagestyle{fancy}
\begin{center}
{\LARGE\color{darkaccent} Алгоритм выбора формулы}
\end{center}
\vspace{13mm}
\begin{center}
\begin{tikzpicture}[>=Latex, node distance=14mm, font=\small,
main/.style={draw=accent,rounded corners=2pt,align=center,text width=71mm,minimum height=10mm,inner sep=4pt},
dec/.style={draw=accent,diamond,aspect=2.65,align=center,text width=42mm,inner sep=2pt,minimum height=13mm},
smallbox/.style={draw=accent,rounded corners=2pt,align=center,text width=43mm,minimum height=12mm,inner sep=3pt},
arr/.style={->,line width=.75pt,draw=darkaccent}]
\node[main] (a) {Прочитайте углы, функции и знаменатели};
\node[main,below=of a] (b) {Вынесите общий множитель, если он есть};
\node[dec,below=17mm of b] (c) {Углы относятся как $2\alpha$ и $\alpha$?};
\node[smallbox,right=27mm of c] (cy) {Используйте формулы двойного угла};
\node[dec,below=24mm of c] (d) {Углы в сумме дают $90^\circ$?};
\node[smallbox,right=27mm of d] (dy) {Используйте формулы дополнительных углов};
\node[main,below=20mm of d] (e) {Выделите кратное $180^\circ$, $360^\circ$ или $2\pi$; проверьте знак};
\node[main,below=17mm of e] (f) {Сведите выражение к одинаковым углам и функциям; сократите ненулевые множители};
\node[main,below=13mm of f] (g) {Вычислите и проверьте знак ответа};
\draw[arr] (a)--(b);
\draw[arr] (b)--(c);
\draw[arr] (c)--node[left,font=\scriptsize]{нет}(d);
\draw[arr] (d)--node[left,font=\scriptsize]{нет}(e);
\draw[arr] (c)--node[above,font=\scriptsize]{да}(cy);
\draw[arr] (d)--node[above,font=\scriptsize]{да}(dy);
\draw[arr] (e)--(f);
\draw[arr] (cy.east)--++(7mm,0) |- ([yshift=2.6mm]f.east);
\draw[arr] (dy.east)--++(4mm,0) |- ([yshift=-2.6mm]f.east);
\draw[arr] (f)--(g);
\end{tikzpicture}
\end{center}
\clearpage

\begin{center}{\LARGE\color{darkaccent} Тренировочный блок}\par\vspace{3mm}{\color{textgray}Путь А \quad $\cdot$ \quad 10 заданий \quad $\cdot$ \quad от простого к сложному}\end{center}
\vspace{7mm}
Во всех упражнениях найдите значение выражения. Преобразования записывайте по шагам; обосновывайте сокращения и проверяйте знаки.
\begin{enumerate}[leftmargin=9mm,itemsep=10pt,topsep=10pt]
\item $\displaystyle \frac{28\cos22^\circ}{\sin68^\circ}+5.$
\item $\displaystyle 35\tg12^\circ\tg78^\circ.$
\item $\displaystyle \frac{23}{\cos^229^\circ+1+\cos^261^\circ}.$
\item $\displaystyle \frac{7\sin104^\circ}{\sin52^\circ\sin38^\circ}.$
\item $\displaystyle \frac{14\sin19^\circ\cos19^\circ}{\sin38^\circ}.$
\item $\displaystyle \frac{18(\sin^223^\circ-\cos^223^\circ)}{\cos46^\circ}.$
\item $\displaystyle \sin\left(-\frac{27\pi}{4}\right)\cos\frac{31\pi}{4}.$
\item $\displaystyle -16\sqrt2\sin(-135^\circ).$
\item $\displaystyle \sqrt{32}\left(\cos^2\frac{13\pi}{8}-\sin^2\frac{13\pi}{8}\right).$
\item $\displaystyle \frac{8\sin130^\circ\cos130^\circ}{\sin100^\circ}.$
\end{enumerate}
\clearpage

\begin{center}{\LARGE\color{darkaccent} Ответы}\par\vspace{4mm}{\color{textgray}Для самостоятельной проверки после решения}\end{center}
\vspace{10mm}
\begin{center}
\renewcommand{\arraystretch}{1.5}
\begin{tabular}{@{}c@{\hspace{28mm}}c@{\hspace{28mm}}c@{\hspace{28mm}}c@{}}
\toprule
№ & Ответ & № & Ответ\\
\midrule
1 & $33$ & 6 & $-18$\\
2 & $35$ & 7 & $-\dfrac12$\\
3 & $\dfrac{23}{2}$ & 8 & $16$\\
4 & $14$ & 9 & $-4$\\
5 & $7$ & 10 & $-4$\\
\bottomrule
\end{tabular}
\end{center}
\vspace{18mm}
\secname{Итоговая памятка}
\renewcommand{\arraystretch}{1.5}
\begin{tabularx}{\linewidth}{@{}p{57mm}X@{}}
\toprule
Сигнал в задаче & Первый полезный шаг\\
\midrule
Сумма углов $90^\circ$ & Выразить один угол через другой и перейти к синусу или косинусу.\\
Один угол вдвое больше & Применить формулы двойного угла.\\
Общий множитель & Вынести за скобки; искать свёртывание формулы.\\
Похожие скобки с противоположными знаками & Вынести минус, затем сокращать.\\
Отрицательный угол & Использовать чётность или нечётность.\\
Большой угол в градусах или радианах & Выделить полный период и табличный остаток.\\
Сокращение дроби & Проверить, что сокращаемый множитель не равен нулю.\\
\bottomrule
\end{tabularx}
\vfill
\begin{center}\small\color{textgray}Лёвин Артём Александрович эксклюзивно для Марка Гукина\end{center}
\end{document}
